%% マニュアル作成要件（2026-10-24-manual-spec.md）§1 の確認に使った最小の文書。
%% jlreq（oneside，版面 37zw）＋ btxdockit で，ltxdockit の環境が和文と一緒に動くことを確かめた。
%% 組版：lualatex docs/proposals/2026-10-24-manual-spike.tex
\documentclass[a4paper,10pt,line_length=37zw,number_of_lines=40,oneside]{jlreq}
\usepackage{luatexja}
\usepackage[haranoaji]{luatexja-preset}
\usepackage{fontspec}
\setmainfont{Libertinus Serif}\setsansfont{Libertinus Sans}
\usepackage{btxdockit}
\renewcommand*{\marglistfont}{\spotcolor\sffamily\small}
\renewcommand*{\optionlistfont}{\spotcolor\sffamily\ttfamily}
\renewcommand*{\ltxsyntaxlabelfont}{\spotcolor\ttfamily}
\hypersetup{colorlinks,linkcolor=spot,urlcolor=spot}
\begin{document}
\section{エントリ型}
\begin{typelist}
\typeitem{book} 単著・共著の書籍。
\reqitem{author, title, date}
\optitem{subtitle, edition, publisher, location}
\end{typelist}
\begin{fieldlist}
\fielditem{origpublisher}{literal} 紙媒体の出版社。電子書籍で使う。
\listitem{location}{literal} 出版地。欧文だけ出す。
\end{fieldlist}
\begin{optionlist}
\boolitem[true]{firstfull} 文献注で，初出の人物をフルネームで出す。
\optitem[false]{doi}{\opt{true}, \opt{false}} 取得日がない文献にも DOI を出す。
\end{optionlist}
\begin{ltxsyntax}
\cmditem{parencite}[prenote][postnote]{key} 全角丸括弧で囲む文献注。
\end{ltxsyntax}
\begin{keymarglist}
\item[translationof] 訳書。原典の項目に添える。
\item[reprintof] 再録・再刊。
\end{keymarglist}
\begin{ltxexample}
@book{aoki2015,
  author = {青木,太郎},
  title  = {日本語の書名},
}
\end{ltxexample}
\end{document}
